\documentclass[12pt,a4paper]{extarticle}
\usepackage[utf8]{inputenc}
\usepackage[T1]{fontenc}
\usepackage[french]{babel}
\usepackage[a4paper,margin=2.2cm]{geometry}
\usepackage{enumitem}
\usepackage{tcolorbox}
\usepackage{longtable}
\usepackage{booktabs}

\setlist[itemize]{nosep}
\setlist[enumerate]{nosep}

\title{\textbf{Correction TD -- Les Fondamentaux de la Gestion d’Incident}\\
BTS SIO}
\author{Maël}
\date{\today}

\begin{document}

\maketitle
\vspace{2cm}
\section*{Exercice 1 -- Classification des situations}

\begin{enumerate}
    \item \textbf{Incident} : Il s'agit d'une interruption non planifiée qui empêche l'utilisateur de réaliser ses missions quotidiennes.
    \item \textbf{Demande} : C'est une requête de service. L'installation n'est pas une panne mais une évolution du poste de travail pour répondre à un besoin métier.
    \item \textbf{Problème} : La répétition de incident indique une cause inconnue. Une analyse est nécessaire pour identifier la source du probleme récurrent.
    \item \textbf{Changement} : Action planifiée et préparée visant à modifier un composant de l'infrastructure réseau pour améliorer ou maintenir le service.
    \item \textbf{Incident} : Interruption d'un service matériel. Le service RH est impacté dans son travail habituel.
    \item \textbf{Demande} : Procédure standard d'accès à une ressource (VPN). Ce n'est pas un incident mais une activation de droits.
    \item \textbf{Problème} : L'aspect systématique et temporel (tous les lundis) suggère un défaut de configuration ou un conflit applicatif latent qu'il faut diagnostiquer.
    \item \textbf{Changement} : Intervention programmée de maintenance corrective ou évolutive sur un système en production.
\end{enumerate}

\vspace{0.5cm}
\clearpage
\section*{Exercice 2 -- Analyse d’impact}

\begin{enumerate}
    \item \textbf{Service concerné :} Direction, Administratif, Technique (Global).\\
          \textbf{Impact :} Critique. arret de l'activité car les données partagées sont inaccessibles. Risque de chômage yechnique immédiat.
    \item \textbf{Service concerné :} Individuel.\\
          \textbf{Impact :} Faible. La gêne est limitée à une personne. Des solutions existent pour contourner temporairement le probleme (imprimer sur un autre poste ou dans un autre service).
    \item \textbf{Service concerné :} Commercial / Marketing. (Dans le cas ou le web n'est pas le coeur de metier de l'entreprise)\\
          \textbf{Impact :} Élevé. Perte de visibilité externe, impossibilité pour les clients de passer commande.
    \item \textbf{Service concerné :} Tous les services.\\
          \textbf{Impact :} Moyen à Élevé. Blocage des flux de communication internes et externes. Retard dans le traitement des dossiers urgents.
\end{enumerate}

\vspace{0.5cm}


\section*{Exercice 3 -- Communication professionnelle}

\begin{enumerate}
    \item \textbf{Réaction :} Garder son calme, faire preuve d'empathie. Il faut laisser l'utilisateur exprimer son probleme sans l'interrompre.
    \item \textbf{Attitude :} Rester neutre, courtois et rassurant. Reformuler le problème pour montrer qu'il a été pris en compte et expliquer les étapes de prise en charge.
    \item \textbf{À éviter :} Entrer dans le conflit, prendre les remarques personnellement, utiliser un jargon technique complexe ou rejeter la faute sur l'utilisateur ou d'autres collègues.
\end{enumerate}

\vspace{0.5cm}

\section*{Exercice 4 -- Typologie des incidents}

\begin{enumerate}
    \item Ordinateur qui ne démarre plus : \textbf{Matériel}
    \item Virus détecté sur un poste : \textbf{Sécurité}
    \item Erreur “404” sur l’intranet : \textbf{Logiciel}
    \item Mot de passe oublié : \textbf{Utilisateur}
    \item Coupure Internet dans tout le bâtiment : \textbf{Réseau}
    \item Écran bleu au démarrage : \textbf{logiciel}
\end{enumerate}

\vspace{0.5cm}


\section*{Exercice 5 -- Cycle de vie simplifié}

\begin{enumerate}
    \item Signalement
    \item Analyse
    \item Résolution
    \item Validation utilisateur
    \item Clôture
\end{enumerate}

\vspace{0.5cm}


\section*{Exercice 6 -- Étude de cas complète}

\begin{enumerate}
    \item \textbf{Nature du cas :} Incident Majeur (incident impactant un grand nombre d'utilisateurs).
    \item \textbf{Impact métier :} Arret des services pour un quart ou un tiers des utilisateurs, empêchant l'accès aux documents de travail.
    \item \textbf{Priorité :} \textbf{Critique}. L'impact est collectif et le service (serveur de fichiers) est vital pour la production.
    \item \textbf{Premières actions :}
    \begin{itemize}
        \item Vérifier l'état physique du serveur (alimentation, connectique).
        \item Tester la connectivité réseau (Ping, services de partage).
        \item Communiquer auprès des utilisateurs via un canal général pour éviter l'engorgement du standard.
    \end{itemize}
    \item \textbf{Importance de la traçabilité :} Elle permet de documenter les solutions trouvées pour la base de connaissances (KEDB), de justifier le temps d'intervention et d'analyser l'incident plus tard s'il se transforme en Problème.
\end{enumerate}

\end{document}